\chapter[Topological Structure of the UOD]{Topological and Geometric Structure of the Universal Optimality Defect}
\section{Introduction}
Let
$F=Q\circ L:\mathcal P^\circ\rightarrow F(\mathcal P^\circ)$
be the global logarithmic embedding of Part III, where
\(\mathcal P^\circ\) denotes the interior of the
probability simplex. The Universal Optimality Defect (UOD) is
$\mathcal D(P,Q)=\|F(P)-F(Q)\|^2.$
Since Part III proves that \(F\) is a global diffeomorphism onto its
image, the topological structure induced by the UOD is obtained by
transporting the Euclidean topology of \(F(\mathcal P^\circ)\).
\section{UOD Balls}
\begin{definition}[UOD Ball]
\label{ch:def:24-1}
For \(P\in\mathcal P^\circ\) and \(r>0\),
$B_r(P)=\{R\in\mathcal P^\circ:\mathcal D(P,R)<r^2\}.$
\end{definition}
\begin{proposition}[UOD Ball Image]
\label{ch:prop:24-2}
$F(B_r(P))=B_r(F(P)),$
where the right-hand side is the Euclidean open ball.
\end{proposition}
\begin{proof}
Since $\mathcal D(P,R)=\|F(P)-F(R)\|^2,$ the claim follows immediately.
\end{proof}
%% Explanatory comment: UOD balls are the inverse images of Euclidean balls under the logarithmic embedding F. This means that every neighborhood in the UOD topology corresponds to a standard Euclidean ball in the image space, establishing the connection between the UOD and Euclidean geometry.
\section{Basis of the Topology}
\begin{theorem}[UOD Balls Form a Basis]
\label{ch:thm:24-3}
The collection of UOD balls forms a basis of the manifold topology on
\(\mathcal P^\circ\).
\end{theorem}
\begin{proof}
Euclidean balls form a basis of the topology of
\(F(\mathcal P^\circ)\). Since \(F\) is a homeomorphism,
their inverse images are precisely the UOD balls. 
\end{proof}
%% Explanatory comment: The UOD balls form a basis of the manifold topology, which means that the topology induced by the UOD is exactly the topology of the posterior manifold. Thus the UOD is a geometrically faithful representation of proximity on the probability simplex.
\section{Convergence}
\begin{theorem}[UOD Convergence]
\label{ch:thm:24-4}
For any sequence \((P_n)\subset\mathcal P^\circ\),
\[ P_n\to P
\quad\Longleftrightarrow\quad
\mathcal D(P_n,P)\to0. \]
\end{theorem}
\begin{proof}
Equivalent to Euclidean convergence of \(F(P_n)\). 
\end{proof}
%% Explanatory comment: Convergence in the UOD is equivalent to convergence in the Euclidean topology of the embedded image. This means that the UOD captures the same notion of sequential convergence as the ambient Euclidean space, making it a practical tool for analysis.
\section{Closed Sets}
\begin{corollary}[Closed Sets]
\label{ch:cor:24-5}
A subset of \(\mathcal P^\circ\) is closed if and only if
it contains the limit of every UOD-convergent sequence.
\end{corollary}
\begin{proof}
Transport the standard sequential characterization of closed subsets of
Euclidean space through the homeomorphism \(F\). 
\end{proof}
%% Explanatory comment: Closed subsets of the posterior simplex can be characterized by UOD convergence, just as in standard Euclidean topology. This allows the application of familiar sequential arguments for closedness in the UOD framework.
\section{Compact Sets}
\begin{theorem}[Compactness]
\label{ch:thm:24-6}
For \(K\subset\mathcal P^\circ\), the following
are equivalent:
\begin{enumerate}
\item \(K\) is compact.
\item Every sequence in \(K\) admits a UOD-convergent subsequence.
\item \(F(K)\) is compact in Euclidean space.
\end{enumerate}
\end{theorem}
\begin{proof}
The equivalence follows because \(F\) is a homeomorphism and Euclidean
spaces are sequentially compact on compact subsets. 
\end{proof}
%% Explanatory comment: Compactness in the UOD topology is equivalent to compactness of the image under F in Euclidean space. This equivalence provides a complete characterization of compact subsets of the posterior simplex, which is essential for existence proofs in optimization.
\section{Local Completeness}
The interior of the simplex is \textbf{not} complete with respect to the
metric induced by the UOD, because sequences may converge to the
boundary, which does not belong to \(\mathcal P^\circ\).
The appropriate completeness statement is therefore local.
\begin{proposition}[Local Completeness]
\label{ch:prop:24-7}
Let \(K\subset\mathcal P^\circ\) be compact.
Every closed subset of \(K\) is complete with respect to the metric
$d(P,Q)=\sqrt{\mathcal D(P,Q)}.$
\end{proposition}
\begin{proof}
Since \(F(K)\) is compact in Euclidean space, every closed subset of
\(F(K)\) is complete. Homeomorphisms preserve Cauchy convergence and
limits for the transported metric, yielding the result. 
\end{proof}
%% Explanatory comment: While the UOD metric is not globally complete (sequences can escape to the boundary), it is complete on compact subsets. This local completeness is sufficient for most analytical purposes: on any bounded region away from the boundary, Cauchy sequences converge within the interior.
\section{Geometric Interpretation}
The UOD introduces no new topology on the posterior manifold. Rather, it
provides an intrinsic realization of the Euclidean topology in
logarithmic coordinates.
All local neighborhoods, convergence properties, compactness arguments
and continuity statements are therefore transported directly from
Euclidean geometry through the logarithmic embedding.
\section{Summary}
The Universal Optimality Defect induces:
\begin{itemize}
\item the intrinsic manifold topology;
\item Euclidean open balls in logarithmic coordinates;
\item the standard notion of sequential convergence;
\item the usual characterization of closed sets;
\item compactness equivalent to Euclidean compactness under \(F\);
\item completeness only on closed subsets of compact regions.
\end{itemize}
The distinction between global incompleteness and local completeness is
essential: the only obstruction to completeness is the excluded boundary
of the probability simplex.
These topological results conclude the intrinsic theory of the UOD. The
following chapter begins its comparison with classical information
geometry, starting from the Fisher metric.
